\section{Introduction}
\label{sec:introduction}

The deployment of autonomous multi-agent systems across critical infrastructure — irrigation networks, financial markets, healthcare logistics, and precision agriculture — has exposed a structural gap in the prevailing architecture of accountability. Modern autonomous agents decompose complex goals into recursive sub-tasks, delegate authority across hierarchies of machine actors, and execute without ongoing human oversight. The prevailing architecture provides no mandatory mechanism for routing exceptional conditions to a human decision-maker when all machine resolution paths have failed. The result is a class of failure modes in which autonomous systems proceed to act without accountability — not because they are adversarial, but because the architecture contains no path back to human consciousness when it is genuinely needed. We call this condition the accountability gap, and the Bailout Protocol is our proposed solution.

\subsection{Why Autonomous Multi-Agent Systems Require a Mandatory Bail-Out Mechanism}

Consider the structure of a deployed autonomous multi-agent system operating under the \bxthree Framework. A Human Accountability Anchor at the root authorizes a \purpose — a bounded objective with explicit constraints. The \bounds engine decomposes that \purpose into sub-tasks, spawning child nodes recursively. Child nodes further decompose and spawn, creating a spawning hierarchy of machine actors that spans from the root anchor to leaf-level actuators. The delegation chain is, in principle, the accountability chain: a human at the top authorizes objectives, intermediate agents translate those objectives into operational decisions, leaf agents execute.

In practice, the chain terminates at the first unresolvable exception. When a leaf agent encounters a condition it cannot resolve with certainty — an out-of-range sensor reading, a safety-envelope violation that its own sandbox simulation cannot rule out, a decision whose implications exceed its provisioned legal or operational authority — it does not escalate to a human. It either fails silently, applies a local heuristic that may or may not be appropriate, or propagates the error upward in a form no human can act on coherently. The exception is absorbed by the machine hierarchy. The human at the top of the chain is never notified. The \purpose layer anchor — the very entity that holds accountability — is structurally bypassed by the autonomous action.

The mandatory bail-out mechanism we propose is not a ``big red button.'' It is an architectural protocol embedded in every node of the spawning hierarchy such that any unresolvable condition at any depth propagates asynchronously to a Human Accountability Anchor, bypassing all intermediate machine actors, with complete diagnostic context from the Forensic Ledger. The protocol is structurally mandatory: it cannot be suppressed by a compromised node, cannot be bypassed by a machine actor in the resolution chain, and cannot be silenced by a network partition without this silence being detected and logged. The structural mandatory property is the essential contribution of this work.

Existing approaches to exception handling in multi-agent systems treat escalation as a design choice to be made at deployment. If the system designer anticipated a condition, they may have included a human review step; if they did not, the condition is handled by the machine hierarchy alone. The Bailout Protocol removes this dependence on anticipation. Any condition that the \bounds engine cannot resolve within its current authority — whether the specific failure mode was anticipated or not — triggers the same mandatory escalation to the \purpose layer. The three trigger conditions that govern this escalation — Capability Boundary, Safety Envelope Violation (predicted), and Accountability Boundary — are defined precisely in Section~\ref{sec:design}, but the critical property is universality: the protocol fires on unresolvable uncertainty, not on a predefined list of error codes.

\subsection{What Happens When Systems Do Not Have One}

The absence of a mandatory bail-out mechanism does not produce neutral outcomes. It produces a specific failure pattern we term \textit{autonomous drift}: the progressive extension of machine authority beyond its intended bounds, without human review, until a failure forces human intervention under suboptimal conditions.

Autonomous drift is enabled by the confidence-threshold model that dominates existing human-in-the-loop frameworks. In this model, a human reviews an agent's action only when the agent's confidence falls below a configured threshold. The model assumes that an agent's self-reported confidence is a reliable signal of whether the action is correct with respect to the user's actual objectives. This assumption fails systematically for large language model--based agents, whose confidence scores are calibrated against the model's training objective (next-token prediction), not against the operational objective (correct action in context). A model can express high confidence while being wrong about what the user intended — and the confidence-threshold model will permit the action to proceed without human review.

The consequences of autonomous drift are compounding. When an agent acts without human review and the action produces an acceptable immediate outcome, the system reinforces the pattern: this agent handled this condition autonomously, no escalation was necessary, the confidence threshold was appropriate. Over successive cycles, machine authority extends into territory that was never authorized, never reviewed, and never recorded in any audit log. The \fact layer — the deterministic enforcement interface at the actuator boundary — executes actions whose upstream authorization has silently expired. The eventual failure is catastrophic in a specific sense: it occurs under conditions the human never reviewed, about whose context they have no diagnostic record, and whose resolution requires intervention in a physical state that has already diverged significantly from the intended state. The human is not a backup safety layer — they are a last-resort responder called in after damage has propagated. This is the inverse of fail-safe design.

The regulatory environment is beginning to respond to this pattern. The EU AI Act and the NIST AI Risk Management Framework both mandate human oversight for high-risk AI systems as a compliance obligation, not a best practice. These mandates presuppose an architecture capable of satisfying them. A system that relies on confidence thresholds and post-hoc audit logs cannot satisfy the mandate because it cannot guarantee that every unresolvable state reached a human before consequential action was taken. The accountability gap is not merely a technical failure — it is a compliance liability.

\subsection{The Core Insight: Accountability Must Have an Architectural Fallback}

The core insight of this paper is that accountability cannot be a property of a system's governance documents — it must be a property of its architecture. Specifically, human accountability must be implemented as a structurally mandatory fallback, not as a configurable threshold or an organizational policy.

This distinction matters because governance documents and architectural constraints have fundamentally different failure modes. A governance document can be updated, overridden, or ignored by a system operator. An architectural constraint cannot be bypassed without modifying the system's structure, and any such modification is itself logged by the Forensic Ledger. When accountability is implemented as a governance document, it exists at the discretion of the system operator. When accountability is implemented as an architectural fallback — as the \bxthree Framework implements it through the Bailout Protocol — it exists independent of any individual operator's choices, including the choices of the humans nominally responsible for the system.

The \bxthree Framework implements this through a three-layer separation of concerns. The \purpose layer owns all Human Accountability Anchors — the named humans with organizational authority who are the terminal recipients of every bailout escalation. The \bounds engine proposes actions within the constraints authorized by the \purpose layer. The \fact layer enforces deterministic physical and operational constraints at the actuator interface. The Bailout Protocol connects these layers: it guarantees that any state the \bounds engine cannot resolve within its current authority — any condition that exceeds the bounds of its provisioned decision space — escalates to a \purpose layer anchor, bypassing all intermediate machine actors, with complete diagnostic context drawn from the Forensic Ledger.

The architectural fallback property has a precise meaning in the \bxthree context. A fallback is not a backup plan that operates when the primary plan fails. It is a mandatory termination condition that the system is structurally required to reach before certain classes of action can proceed. In the \bxthree architecture, the termination condition is a Human Accountability Anchor, and the mandatory property is enforced by two structural constraints: the immutability of parent pointers in the spawning hierarchy, and the bypass mechanism of the Bailout Protocol itself. No agent can modify its own parent pointer. No machine actor in the resolution chain can redirect an escalation away from its designated Human Accountability Anchor without this redirection being detected and logged as an anomaly. The accountability chain does not rely on the good judgment of system operators. It is structurally enforced.

This is the guarantee that no existing multi-agent framework provides. Existing systems implement exception handling as a design choice, confidence thresholds as a configurable parameter, and human oversight as an organizational policy. The \bxthree Framework implements human oversight as a mathematical consequence of immutable parent pointers, depth-bounded spawning, and the mandatory termination property of the Bailout Protocol.

\subsection{Paper Roadmap}

The remainder of this paper is organized as follows. Section~\ref{sec:background} reviews the foundational threads converging in the Bailout Protocol: the principle of separation of concerns as an accountability architecture, the problem of accountability in multi-agent AI systems, human-in-the-loop design and its structural limitations, and fail-safe design in safety-critical systems. Section~\ref{sec:design} presents the formal design of the Bailout Protocol, including its three triggering conditions, escalation state machine, and the bypass mechanism that severs all machine actors upon activation. Section~\ref{sec:forensic-ledger} describes the Forensic Ledger as the audit substrate on which the Bailout Protocol operates, including the 9-Plane Data Action Protocol structure and the hash-chained append-only architecture that guarantees log integrity. Section~\ref{sec:properties} establishes the protocol's correctness properties: bypass soundness, no suppression by compromised nodes, and deterministic resolution at the Human Accountability Anchor. Section~\ref{sec:relationship} situates the Bailout Protocol within the broader \bxthree Framework, including its relationship to the Training Wheels Protocol and the Human Accountability Anchor specification. Section~\ref{sec:limitations} acknowledges the protocol's limitations — human response latency under high-frequency events, adversarial conditions, deployment depth scalability, and exploitation attack vectors — and proposes directions for future work. Section~\ref{sec:conclusion} concludes with a discussion of the protocol's implications for mandatory human accountability in autonomous systems and its relationship to emerging regulatory frameworks.